\chapter{Solution Design}
\label{ch:design}

\vcyclephase{Functional Design + Detailed Design (Middle of V)}

\section{System Architecture}
\label{sec:sys-arch}

The SSMS follows a three-tier architecture: the desktop application (presentation and business logic), the MQTT broker (communication), and the ESP32 with peripherals (physical control).

\begin{figure}[H]
  \centering
  \begin{tikzpicture}[
      node distance=2cm,
      block/.style={rectangle, draw, minimum width=3cm, minimum height=1.2cm, align=center, fill=ssmsBlue!10, thick, font=\small},
      arrow/.style={-{Latex[length=2mm]}, thick},
      label/.style={font=\tiny, above}
    ]
    \node[block] (app) {Desktop App\\(Python/Tkinter)};
    \node[block] (broker) [right=of app] {MQTT Broker\\(Mosquitto)};
    \node[block] (esp) [right=of broker] {ESP32\\(MicroPython)};
    \node[block] (level) [right=of esp] {Level Shifter\\(74HCT244/245/573)};
    \node[block] (leds) [right=of level] {WS2813\\LED Strip};

    \draw[arrow, <->] (app) -- node[label] {MQTT} (broker);
    \draw[arrow, <->] (broker) -- node[label] {MQTT} (esp);
    \draw[arrow, ->] (esp) -- node[label] {3.3V GPIO} (level);
    \draw[arrow, ->] (level) -- node[label] {5V Data} (leds);
  \end{tikzpicture}
  \caption{System architecture block diagram.}
  \label{fig:sys-arch}
\end{figure}

\subsection{Power Architecture}
\label{sec:power-arch}

Two separate power rails are used to ensure stable operation: 3.3\,V for the ESP32 via the HLK-PM01 AC/DC module, and 5\,V for the LED strip via a dedicated 5\,V/5\,A PSU.

\begin{figure}[H]
  \centering
  \begin{tikzpicture}[
      block/.style={rectangle, draw, minimum width=2.5cm, minimum height=1cm, align=center, font=\small, thick},
      arrow/.style={-{Latex[length=2mm]}, thick}
    ]
    \node[block] (mains) at (0,0) {Mains\\(230 VAC)};
    \node[block] (hlk) [right=of mains] {HLK-PM01\\(5V, 600mA)};
    \node[block] (esp) [right=of hlk] {ESP32\\(3.3V)};
    \node[block] (psu) [below=of mains] {5V/5A\\PSU};
    \node[block] (leds) [right=of psu] {WS2813\\LEDs (5V)};

    \draw[arrow, ->] (mains) -- (hlk);
    \draw[arrow, ->] (hlk) -- (esp);
    \draw[arrow, ->] (mains) -- (psu);
    \draw[arrow, ->] (psu) -- (leds);
  \end{tikzpicture}
  \caption{Power architecture diagram.}
  \label{fig:power-arch}
\end{figure}

\section{Hardware Design}
\label{sec:hw-design}

\subsection{ESP32-WROOM-32D Pin Assignment}
\label{sec:pin-assignment}

\begin{table}[H]
  \centering
  \caption{ESP32 pin assignment.}
  \label{tab:pins}
  \begin{tabular}{p{2cm} p{3cm} p{6cm}}
    \toprule
    \textbf{Pin} & \textbf{Function} & \textbf{Connection} \\
    \midrule
    GPIO18 & LED data out & To 74HCT245 input (level shift to 5\,V) \\
    GPIO19 & MQTT status LED & On-board indicator (optional) \\
    GPIO21 & I2C SDA & Reserved for future expansion \\
    GPIO22 & I2C SCL & Reserved for future expansion \\
    GPIO25 & Debug UART TX & Reserved for debugging \\
    GPIO26 & Debug UART RX & Reserved for debugging \\
    EN & Reset & Pull-up with 10\,k$\Omega$ to 3.3\,V \\
    \bottomrule
  \end{tabular}
\end{table}

\subsection{Level-Shifting Circuit}
\label{sec:level-shift}

The WS2813 LED strip requires a 5\,V data signal, but the ESP32 GPIO outputs operate at 3.3\,V. Three 74HCT-series ICs are employed:

\begin{itemize}
  \item \textbf{SN74HCT244:} Octal buffer/line driver. Used to buffer and level-shift the data signal from 3.3\,V to 5\,V.
  \item \textbf{SN74HCT245:} Octal bus transceiver. Used for bidirectional level shifting if needed; in this design, it operates in the A-to-B direction.
  \item \textbf{SN74HCT573:} Octal transparent D-type latch. Provides additional drive strength and can optionally latch data for parallel output expansion.
\end{itemize}

The 74HCT series is specifically chosen because its VIH threshold is 2.0\,V (compatible with 3.3\,V logic)\cite{sn74hct573_ds}, unlike the 74HC series which has a VIH of approximately 3.5\,V. This ensures reliable switching with the ESP32's 3.3\,V GPIO output.

\subsection{WS2813 Wiring Topologies}
\label{sec:wiring}

The 36 bins can be wired in three different topologies:

\begin{table}[H]
  \centering
  \caption{LED strip wiring topology comparison.}
  \label{tab:topology}
  \begin{tabular}{p{2.5cm} p{3.5cm} p{3.5cm} p{3.5cm}}
    \toprule
    \textbf{Property} & \textbf{Snake (1 GPIO)} & \textbf{Column (4 GPIO)} & \textbf{Row (9 GPIO)} \\
    \midrule
    GPIOs used & 1 & 4 & 9 \\
    Wiring complexity & Simple & Medium & Complex \\
    Individual bin control & Yes & Yes & Yes \\
    Max refresh rate & 216 pixels @ 800\,kHz & 54 pixels per GPIO & 24 pixels per GPIO \\
    Fault isolation & None (single strip) & Per column & Per row \\
    \bottomrule
  \end{tabular}
\end{table}

\subsection{Protection Components}
\label{sec:protection}

The following protection components are included on the PCB:

\begin{itemize}
  \item \textbf{SS54 / SS14:} Schottky flyback diodes across the LED power rail to protect against inductive kickback from long LED strip cables.
  \item \textbf{SMBJ5.0A:} TVS diode for overvoltage protection on the 5\,V supply rail.
  \item \textbf{33--100\,$\Omega$ series resistor:} Placed on the LED data line near the source to dampen ringing and match impedance.
\end{itemize}

\subsection{KiCad Schematic}
\label{sec:schematic}

The full KiCad schematic is included in Annex E. \textit{[Placeholder: Insert KiCad schematic export PDF/PNG here]}

\subsection{Bill of Materials}
\label{sec:bom}

\begin{longtable}{p{3.5cm} p{3.5cm} p{1.5cm} p{2cm} p{2cm}}
  \caption{Bill of Materials.}\label{tab:bom}\\
  \toprule
  \textbf{Component} & \textbf{Part Number} & \textbf{Qty} & \textbf{Unit Cost (USD)} & \textbf{Total (USD)} \\
  \midrule
  \endfirsthead
  \multicolumn{5}{c}{\tablename\ \thetable\ --- Continued from previous page} \\
  \toprule
  \textbf{Component} & \textbf{Part Number} & \textbf{Qty} & \textbf{Unit Cost} & \textbf{Total} \\
  \midrule
  \endhead
  \bottomrule
  \endfoot
  MCU module & ESP32-WROOM-32D & 1 & 5.50 & 5.50 \\
  RGB LED strip & WS2813 (per metre) & 12 & 8.00 & 96.00 \\
  Level shifter & SN74HCT244 & 1 & 0.85 & 0.85 \\
  Bus transceiver & SN74HCT245 & 1 & 0.90 & 0.90 \\
  D-type latch & SN74HCT573 & 1 & 0.80 & 0.80 \\
  AC/DC module & HLK-PM01 & 1 & 3.50 & 3.50 \\
  LED PSU & 5V/5A switching & 1 & 12.00 & 12.00 \\
  Schottky diode & SS54 & 2 & 0.30 & 0.60 \\
  TVS diode & SMBJ5.0A & 1 & 0.25 & 0.25 \\
  Resistor kit & Various (0603) & 1 & 2.00 & 2.00 \\
  Capacitor kit & Various (0603) & 1 & 3.00 & 3.00 \\
  PCB fabrication & Custom 2-layer & 1 & 15.00 & 15.00 \\
  Connectors & JST-XH 2-pin & 10 & 0.50 & 5.00 \\
  \midrule
  \multicolumn{4}{r}{\textbf{Total}} & \textbf{145.40} \\
\end{longtable}

\section{Firmware Design}
\label{sec:fw-design}

\subsection{MicroPython Architecture}
\label{sec:fw-arch}

The firmware is structured into three modules:

\begin{itemize}
  \item \texttt{main.py}: Initialisation, WiFi connection, MQTT client setup, main loop.
  \item \texttt{mqtt\_handler.py}: MQTT callback functions, message parsing, command dispatch.
  \item \texttt{led\_controller.py}: NeoPixel initialisation, LED state machine, blink timer.
\end{itemize}

\subsection{LED State Machine}
\label{sec:state-machine}

The firmware implements a finite state machine with five states:

\begin{figure}[H]
  \centering
  \begin{tikzpicture}[
      node distance=2.5cm,
      state/.style={circle, draw, minimum size=1.2cm, align=center, font=\small\bfseries, thick},
      arrow/.style={-{Latex[length=2mm]}, thick},
      label/.style={font=\tiny, sloped}
    ]
    \node[state, fill=ssmsGray] (idle) {IDLE};
    \node[state, fill=ssmsRed!20] (pending) [right=of idle] {PENDING};
    \node[state, fill=ssmsBlue!20] (active) [right=of pending] {ACTIVE};
    \node[state, fill=ssmsOrange!20] (scanning) [below=of active] {SCANNING};
    \node[state, fill=ssmsGreen!20] (complete) [below=of pending] {COMPLETE};

    \draw[arrow, ->] (idle) -- node[label, above] {LED\_ON} (pending);
    \draw[arrow, ->] (pending) -- node[label, above] {scan} (active);
    \draw[arrow, ->] (active) -- node[label, right] {scanning} (scanning);
    \draw[arrow, ->] (scanning) -- node[label, below] {confirm} (complete);
    \draw[arrow, ->] (complete) -- node[label, left] {2\,s timeout} (idle);
    \draw[arrow, ->] (active) to[out=135,in=45] node[label, above] {LED\_OFF} (idle);
    \draw[arrow, ->] (scanning) to[out=180,in=0] node[label, above] {LED\_OFF} (idle);
  \end{tikzpicture}
  \caption{LED state machine diagram.}
  \label{fig:state-machine}
\end{figure}

\subsection{MQTT Message Parsing}
\label{sec:mqtt-parsing}

Commands received on the \texttt{ssms/command} topic are parsed and dispatched to the LED controller:

\begin{figure}[H]
  \centering
  \begin{tikzpicture}[
      node distance=1.8cm,
      block/.style={rectangle, draw, minimum width=2.5cm, minimum height=0.8cm, align=center, font=\footnotesize, thick},
      decision/.style={diamond, draw, aspect=2, minimum width=1.5cm, minimum height=1cm, align=center, font=\footnotesize, thick},
      arrow/.style={-{Latex[length=2mm]}, thick}
    ]
    \node[block] (start) {MQTT msg\\received};
    \node[decision] (parse) [below=of start] {Parse\\command};
    \node[block] (ledon) [below left=of parse] {LED\_ON};
    \node[block] (ledoff) [below=of parse] {LED\_OFF};
    \node[block] (allon) [below right=of parse] {ALL\_ON/OFF};
    \node[block] (blink) [right=of parse] {BLINK};
    \node[block] (exec) [below=of ledoff] {Execute\\LED action};
    \node[block] (publish) [below=of exec] {Publish\\status};

    \draw[arrow] (start) -- (parse);
    \draw[arrow] (parse) -- node[label, above] {LED\_ON} (ledon);
    \draw[arrow] (parse) -- node[label, right] {LED\_OFF} (ledoff);
    \draw[arrow] (parse) -- node[label, above] {ALL\_*} (allon);
    \draw[arrow] (parse) -- node[label, above] {BLINK} (blink);
    \draw[arrow] (ledon) |- (exec);
    \draw[arrow] (ledoff) -- (exec);
    \draw[arrow] (allon) |- (exec);
    \draw[arrow] (blink) |- (exec);
    \draw[arrow] (exec) -- (publish);
  \end{tikzpicture}
  \caption{MQTT message parsing flowchart.}
  \label{fig:mqtt-flow}
\end{figure}

\subsection{NeoPixel Addressing Formula}
\label{sec:neopixel-addr}

Each bin is assigned a unique group of 6 consecutive LEDs. The pixel index for a given bin and LED offset is calculated as:

\begin{equation}
  \text{pixel\_index} = \text{bin\_number} \times \text{LEDS\_PER\_BIN} + \text{offset}
  \label{eq:neopixel}
\end{equation}

Where $\text{LEDS\_PER\_BIN} = 6$ and $\text{offset} \in \{0, 1, 2, 3, 4, 5\}$.

\section{Desktop Application Design}
\label{sec:app-design}

\subsection{Architecture: MVC Pattern}
\label{sec:mvc}

The desktop application follows the Model-View-Controller (MVC) pattern with four modules:

\begin{itemize}
  \item \texttt{csv\_mgr.py} (Model): Handles CSV data loading and validation.
  \item \texttt{history\_mgr.py} (Model): Manages pick history logging and retrieval.
  \item \texttt{mqtt\_client.py} (Controller): Wraps the paho-mqtt client with application-specific methods.
  \item \texttt{ui\_main.py} (View): Tkinter GUI layout, widgets, and event binding.
\end{itemize}

\subsection{Workflow Diagram}
\label{sec:workflow}

\begin{figure}[H]
  \centering
  \begin{tikzpicture}[
      node distance=2cm,
      block/.style={rectangle, draw, minimum width=2.5cm, minimum height=0.8cm, align=center, font=\footnotesize, thick},
      arrow/.style={-{Latex[length=2mm]}, thick}
    ]
    \node[block] (scanbadge) {Scan\\badge};
    \node[block] (verify) [right=of scanbadge] {Verify\\operator};
    \node[block] (batch) [right=of verify] {Load\\batch};
    \node[block] (binscan) [right=of batch] {Scan\\bin};
    \node[block] (partscan) [right=of binscan] {Scan\\part};
    \node[block] (complete) [below=of partscan] {Completion\\+ log};

    \draw[arrow, ->] (scanbadge) -- (verify);
    \draw[arrow, ->] (verify) -- (batch);
    \draw[arrow, ->] (batch) -- (binscan);
    \draw[arrow, ->] (binscan) -- (partscan);
    \draw[arrow, ->] (partscan) -- (complete);
  \end{tikzpicture}
  \caption{Desktop application workflow.}
  \label{fig:workflow}
\end{figure}

\subsection{GUI Wireframe}
\label{sec:gui}

The main application window is divided into three sections:

\begin{enumerate}
  \item \textbf{Header bar:} Operator name, current batch ID, timer.
  \item \textbf{Main panel:} Table view of current picks showing bin location, item code, quantity, status.
  \item \textbf{Status bar:} MQTT connection status, last action, error messages.
\end{enumerate}

\subsection{CSV Data Model}
\label{sec:csv}

The CSV file used for batch loading has the following columns:

\begin{table}[H]
  \centering
  \caption{CSV data schema.}
  \label{tab:csv}
  \begin{tabular}{p{3cm} p{2.5cm} p{6cm}}
    \toprule
    \textbf{Column} & \textbf{Type} & \textbf{Description} \\
    \midrule
    bin\_id & Integer & Bin number (0--35) \\
    item\_code & String & Product SKU or barcode \\
    item\_description & String & Human-readable product name \\
    quantity & Integer & Quantity to pick \\
    operation & String & ``PULL'' or ``PUT'' \\
    \bottomrule
  \end{tabular}
\end{table}

\section{LTspice Simulation Design}
\label{sec:ltspice-design}

\subsection{Simulation Objective}
\label{sec:ltspice-objective}

The LTspice simulation was used to verify the level-shifter circuit performance. Specifically, it confirmed:

\begin{itemize}
  \item The 3.3\,V GPIO output is correctly shifted to a 5\,V signal at the WS2813 data input.
  \item The signal integrity is maintained at the WS2813 communication frequency of 800\,kHz.
  \item The pull-up resistor value on the data line is correctly sized.
\end{itemize}

\subsection{Methodology}
\label{sec:ltspice-method}

The simulation model included the ESP32 GPIO output (modelled as a 3.3\,V pulse source with 5\,ns rise/fall time), the 74HCT245 buffer, and a simplified load representing the WS2813 input capacitance. The output was measured at the WS2813 data pin.

\section{Proteus Simulation Design}
\label{sec:proteus-design}

\subsection{Simulation Objective}
\label{sec:proteus-objective}

The Proteus simulation was used to verify the digital logic behaviour of the ESP32 communicating with the 74HCT573 latch.

\subsection{Known Limitation}
\label{sec:proteus-limitation}

During simulation, it was observed that the SN74HCT573 model in Proteus behaves as a standard 74HC device (VIH $\approx$ 3.5\,V), rather than a true 74HCT device (VIH = 2.0\,V). As a result, the 3.3\,V input from the ESP32 produces a 0\,V output in simulation. This is a \textbf{simulator artifact} due to an inaccurate model --- the Texas Instruments SN74HCT573 datasheet\cite{sn74hct573_ds} confirms a maximum VIH of 2.0\,V, guaranteeing correct operation at 3.3\,V in the real circuit.
